% Sección 4.3: La Integral Definida

Ejercicios resueltos sobre la definición formal de la integral definida mediante el límite de sumas de Riemann, integrabilidad y propiedades fundamentales.

\section*{SECCIÓN 4.3 \quad Ejercicios}

\noindent \textit{En los Ejercicios 1 a 6 se da una función $f$, un intervalo, puntos de división que definen una partición $P$ y puntos $x_i^*$ del $i$-ésimo subintervalo. (a) Calcule $\|P\|$. (b) Obtenga la suma de Riemann (4.10).}

\begin{enumerate}
    \subimport{ejercicio_01/}{ejercicio.tex}
    \subimport{ejercicio_02/}{ejercicio.tex}
    \subimport{ejercicio_03/}{ejercicio.tex}
    \subimport{ejercicio_04/}{ejercicio.tex}
    \subimport{ejercicio_05/}{ejercicio.tex}
    \subimport{ejercicio_06/}{ejercicio.tex}
\end{enumerate}

\vspace{1em}
\noindent \textit{En los Ejercicios 7 a 10, aplique la Regla del Punto Medio con el valor dado en $n$ para aproximar cada integral. Redondee sus respuestas hasta cuatro cifras decimales.}

\begin{enumerate}
    \setcounter{enumi}{6}
    \begin{minipage}{0.45\textwidth}
    \subimport{ejercicio_07/}{ejercicio.tex}
    \subimport{ejercicio_08/}{ejercicio.tex}
    \end{minipage}%
    \begin{minipage}{0.45\textwidth}
    \subimport{ejercicio_09/}{ejercicio.tex}
    \subimport{ejercicio_10/}{ejercicio.tex}
    \end{minipage}
\end{enumerate}

\vspace{1em}
\noindent \textit{Aplique el Teorema 4.12 para evaluar la integral de los Ejercicios 11 a 22.}

\begin{enumerate}
    \setcounter{enumi}{10}
    \begin{minipage}{0.45\textwidth}
    \subimport{ejercicio_11/}{ejercicio.tex}
    \subimport{ejercicio_12/}{ejercicio.tex}
    \subimport{ejercicio_13/}{ejercicio.tex}
    \subimport{ejercicio_14/}{ejercicio.tex}
    \subimport{ejercicio_15/}{ejercicio.tex}
    \subimport{ejercicio_16/}{ejercicio.tex}
    \end{minipage}%
    \begin{minipage}{0.45\textwidth}
    \subimport{ejercicio_17/}{ejercicio.tex}
    \subimport{ejercicio_18/}{ejercicio.tex}
    \subimport{ejercicio_19/}{ejercicio.tex}
    \subimport{ejercicio_20/}{ejercicio.tex}
    \subimport{ejercicio_21/}{ejercicio.tex}
    \subimport{ejercicio_22/}{ejercicio.tex}
    \end{minipage}
\end{enumerate}

\vspace{1em}
\begin{enumerate}
    \setcounter{enumi}{22}
    \subimport{ejercicio_23/}{ejercicio.tex}
    \subimport{ejercicio_24/}{ejercicio.tex}
\end{enumerate}

\vspace{1em}
\noindent \textit{Exprese los límites indicados en los Ejercicios 25 a 28 como una integral definida en el intervalo dado.}

\begin{enumerate}
    \setcounter{enumi}{24}
    \subimport{ejercicio_25/}{ejercicio.tex}
    \subimport{ejercicio_26/}{ejercicio.tex}
    \subimport{ejercicio_27/}{ejercicio.tex}
    \subimport{ejercicio_28/}{ejercicio.tex}
\end{enumerate}

\vspace{1em}
\noindent \textit{Exprese los límites indicados en los Ejercicios 29 a 31 como una integral definida.}

\begin{enumerate}
    \setcounter{enumi}{28}
    \subimport{ejercicio_29/}{ejercicio.tex}
    \subimport{ejercicio_30/}{ejercicio.tex}
    \subimport{ejercicio_31/}{ejercicio.tex}
\end{enumerate}

\vspace{1em}
\begin{enumerate}
    \setcounter{enumi}{31}
    \subimport{ejercicio_32/}{ejercicio.tex}
    \subimport{ejercicio_33/}{ejercicio.tex}
    \subimport{ejercicio_34/}{ejercicio.tex}
    \subimport{ejercicio_35/}{ejercicio.tex}
    \subimport{ejercicio_36/}{ejercicio.tex}
    \subimport{ejercicio_37/}{ejercicio.tex}
    \subimport{ejercicio_38/}{ejercicio.tex}
    \subimport{ejercicio_39/}{ejercicio.tex}
\end{enumerate}
